\chapter{Charged Particles, Field Transformations, and Relativistic Waves}

Last time we began asking how a charged particle moves in electric and
magnetic fields once the principles of relativity are imposed. The familiar
Lorentz force law already contains the clue: its magnetic term depends on the
particle's velocity, so different observers cannot agree separately about
what is electric and what is magnetic. Our task is to formulate the law so
that every inertial observer uses the same equation, even though the fields
and the particle velocity change from frame to frame.

Unless stated otherwise, we use units in which \(c=1\) and the metric
signature \((-,-,-,+)\). Spacetime components are ordered with the three
spatial entries first:
\begin{equation}
  \mu=(1,2,3,0)=(x,y,z,t).
\end{equation}
Latin indices \(m,n\) run only over the three spatial directions.

\section{What survives of Newton's force law}

In Newtonian mechanics,
\begin{equation}
  \boldsymbol F=m\boldsymbol a
\end{equation}
has the same form in every inertial frame because Galilean observers agree
on acceleration. They also assign the same mass and force to the particle.
Once the kinematics changes, there is no reason to expect this particular
relation between force and ordinary acceleration to survive unchanged.

Electromagnetism and gravitation already warn us that relativistic
extensions need not resemble one another. Coulomb's law and Newton's law of
gravity are both inverse-square laws, but gravitation cannot be incorporated
naturally into special relativity without a deeper change in spacetime
geometry. Electromagnetism can. Indeed, magnetism is itself a relativistic
effect: when factors of \(c\) are restored, a magnetic force carries a factor
of velocity divided by \(c\), and it disappears in the formal limit
\(c\to\infty\).

The durable part of Newton's law is the statement that force is the rate of
change of momentum. For constant rest mass,
\begin{equation}
  \boldsymbol F
  =m\frac{d\boldsymbol v}{dt}
  =\frac{d}{dt}(m\boldsymbol v)
  =\frac{d\boldsymbol p}{dt}.
\end{equation}
Relativity will alter the relation between momentum and velocity, but it
will preserve this way of organizing the equation.

\begin{figure}[H]
\centering
\includegraphics[width=0.64\textwidth]{lecture_07_figure_01.png}
\caption{Force rewritten as \(d\boldsymbol p/dt\) while the electric part of
the Lorentz law is being placed on the board.}
\lectureframe{00:08:20}
\end{figure}

In units \(c=1\), the pre-relativistic Lorentz force law is
\begin{equation}
  \frac{d\boldsymbol p}{dt}
  =q\bigl(\boldsymbol E+\boldsymbol v\times\boldsymbol B\bigr).
  \label{eq:lorentz-three}
\end{equation}
Both fields may depend on position and time. The nonrelativistic limit is
now the limit \(v\ll 1\), not a limit in which the numerical constant \(c=1\)
somehow changes. The electric term in \eqref{eq:lorentz-three} is already an
ordinary vector. To understand the magnetic term in a form that can be
carried into spacetime, we must look again at the cross product.

\section{The cross product as an antisymmetric tensor}

Take two three-dimensional vectors \(\boldsymbol A\) and \(\boldsymbol B\).
Their outer product
\begin{equation}
  T_{mn}=A_mB_n
\end{equation}
is a rank-two tensor with nine components. It has no symmetry in general.
Antisymmetrizing gives
\begin{equation}
  K_{mn}=A_mB_n-A_nB_m,
  \qquad K_{mn}=-K_{nm}.
  \label{eq:antisym-product}
\end{equation}
Its diagonal entries vanish, while every entry below the diagonal is fixed
by the corresponding entry above it. A \(3\times3\) antisymmetric tensor
therefore has only three independent components.

That count is suggestive but not sufficient by itself. Three quantities do
not automatically transform as a vector. In three dimensions, however, the
three independent entries of \eqref{eq:antisym-product} can be identified
with the components of the cross product
\(\boldsymbol C=\boldsymbol A\times\boldsymbol B\). The board first makes a
temporary placement of \(C_2\) in the upper-right position:

\begin{figure}[H]
\centering
\includegraphics[width=0.75\textwidth]{lecture_07_figure_02.png}
\caption{The antisymmetric product while its three independent entries are
being matched to the components of \(\boldsymbol C\). The upper-right label
is still provisional here.}
\lectureframe{00:17:33}
\end{figure}

Cycling \(1\to2\to3\to1\) gives the component formulas
\begin{align}
  C_3 &= A_1B_2-A_2B_1, \\
  C_1 &= A_2B_3-A_3B_2, \\
  C_2 &= A_3B_1-A_1B_3.
  \label{eq:cross-components}
\end{align}
The last line shows that the upper-right entry is \(-C_2\), not \(C_2\).
The correction is part of the argument rather than a detail to hide.

\begin{figure}[H]
\centering
\includegraphics[width=0.62\textwidth]{lecture_07_figure_03.png}
\caption{The completed cyclic component identities that fix the temporary
sign in the antisymmetric array.}
\lectureframe{00:21:54}
\end{figure}

The resulting matrix is
\begin{equation}
  K_{mn}=
  \begin{pmatrix}
    0 & C_3 & -C_2 \\
    -C_3 & 0 & C_1 \\
    C_2 & -C_1 & 0
  \end{pmatrix}.
  \label{eq:cross-matrix}
\end{equation}

There is a useful geometry behind this notation. A component with indices
\(1,2\) belongs naturally to the \(xy\)-plane and is identified with a
vector component along the perpendicular \(z\)-axis. One may describe the
same oriented object either by a plane or by an arrow normal to that plane.
The arrow description is special to three dimensions; in higher dimensions
an antisymmetric rank-two tensor still exists, but it no longer has the same
number of components as a vector. Thus the tensor, not the perpendicular
arrow, is the form of the cross product that generalizes.

Equivalently, after choosing an orientation,
\begin{equation}
  C_k=\frac12\epsilon_{kmn}K_{mn}.
\end{equation}
The need to choose an orientation is the source of the right-hand rule and
of the familiar distinction between polar and axial vectors.
\editorialnote{The Levi--Civita formula and the axial-vector terminology are
standard notation for the plane-versus-normal geometry developed explicitly
in the lecture.}

\section{The magnetic field in tensor form}

The same dual description applies to the magnetic field. We define its
antisymmetric spatial tensor by
\begin{equation}
  B_{mn}=
  \begin{pmatrix}
    0 & -B_3 & B_2 \\
    B_3 & 0 & -B_1 \\
    -B_2 & B_1 & 0
  \end{pmatrix}.
  \label{eq:magnetic-tensor}
\end{equation}
The extra overall sign relative to \eqref{eq:cross-matrix} is a convention.
Once it has been chosen, it must be used consistently; changing the sign of
all three vector components would describe the same tensor-to-vector
identification with the opposite orientation convention.

\begin{figure}[H]
\centering
\includegraphics[width=0.72\textwidth]{lecture_07_figure_04.png}
\caption{The magnetic field written both as an antisymmetric spatial array
and in terms of its three vector components.}
\lectureframe{00:24:40}
\end{figure}

Now take the third component of the cross product:
\begin{align}
  (\boldsymbol v\times\boldsymbol B)_3
  &=v_1B_2-v_2B_1 \\
  &=v_1B_{13}+v_2B_{23}+v_3B_{33}.
\end{align}
The last term may be added because \(B_{33}=0\). Cycling the free component
gives the compact rule
\begin{equation}
  (\boldsymbol v\times\boldsymbol B)_m=v_nB_{nm},
  \label{eq:cross-contraction}
\end{equation}
where the repeated index \(n\) is summed. The magnetic force is therefore
\begin{equation}
  F_m^{(B)}=qv_nB_{nm}.
\end{equation}

\begin{figure}[H]
\centering
\includegraphics[width=0.78\textwidth]{lecture_07_figure_05.png}
\caption{The third component of \(\boldsymbol v\times\boldsymbol B\) and its
general expression as a contraction with \(B_{nm}\).}
\lectureframe{00:28:30}
\end{figure}

\begin{classroomqa}
\audiencequestion{Does the free index \(m\) label a separate magnetic
component of the force?}
\lecturerresponse{No. It labels whichever spatial component of the force is
being calculated. The electric field has simply been set to zero for the
present argument, and every index here runs only over the three spatial
directions.}
\end{classroomqa}

The tensor notation makes the no-work property of a magnetic field nearly
automatic. Differentiate the squared speed:
\begin{align}
  \frac{d}{dt}(v_mv_m)
  &=2\dot v_m v_m \\
  &\propto 2(\boldsymbol v\times\boldsymbol B)_m v_m \\
  &=2v_nB_{nm}v_m.
  \label{eq:speed-derivative}
\end{align}
Terms with \(m=n\) vanish because the diagonal of \(B_{nm}\) is zero. Every
off-diagonal term has a partner with \(m\) and \(n\) interchanged; the two
cancel because \(v_nv_m\) is symmetric while \(B_{nm}\) is antisymmetric.
Hence
\begin{equation}
  \frac{d}{dt}(v_mv_m)=0.
\end{equation}
A pure magnetic force bends a trajectory without changing its speed or
kinetic energy.

\begin{figure}[H]
\centering
\includegraphics[width=0.68\textwidth]{lecture_07_figure_06.png}
\caption{The paired terms in \(v_nB_{nm}v_m\) cancel because the magnetic
tensor changes sign when its indices are interchanged.}
\lectureframe{00:35:40}
\end{figure}

This proof is intentionally more elaborate than saying that a cross product
is perpendicular to its first factor. The latter statement is peculiar to
the three-dimensional vector representation. The contraction of a symmetric
pair with an antisymmetric tensor is the argument that will survive in four
dimensions.

\begin{classroomqa}
\audiencequestion{Must the magnetic field be time independent for the speed
to remain constant?}
\lecturerresponse{The algebra above never differentiates the field and does
not require \(\partial_t\boldsymbol B=0\). A magnetic term by itself still
does no work. In electrodynamics, however, a changing magnetic field produces
an electric field, and that electric field can change the particle's speed.}
\end{classroomqa}

\section{Why electric and magnetic fields must mix}

The mixing of \(\boldsymbol E\) and \(\boldsymbol B\) can be seen before any
four-dimensional calculation. Consider slow motion, so that ordinary
Newtonian statements about acceleration are adequate. Let a magnetic field
point along the \(y\)-axis and let a charged particle move along the
\(x\)-axis. The magnetic force points along the \(z\)-axis:
\begin{equation}
  \boldsymbol F=q\boldsymbol v\times\boldsymbol B.
\end{equation}

\begin{figure}[H]
\centering
\includegraphics[width=0.62\textwidth]{lecture_07_figure_07.png}
\caption{The board sketch for a particle moving across a magnetic field;
the three directions are used to compare the laboratory and co-moving
descriptions.}
\lectureframe{00:42:20}
\end{figure}

The geometry may be redrawn without the perspective of the classroom:
\begin{figure}[H]
\centering
\begin{tikzpicture}[scale=1.05,>=Stealth]
  \draw[->] (0,0) -- (3.8,0) node[right] {\(x\)};
  \draw[->] (0,0) -- (-1.9,2.4) node[above left] {\(y\)};
  \draw[->] (0,0) -- (0,3.1) node[above] {\(z\)};
  \draw[->,very thick] (0,0) -- (2.7,0)
    node[midway,below] {\(\boldsymbol v\)};
  \draw[->,very thick] (0,0) -- (-1.35,1.7)
    node[midway,above left] {\(\boldsymbol B\)};
  \draw[->,very thick] (0,0) -- (0,2.35)
    node[midway,right] {\(\boldsymbol F,\ \boldsymbol E'\)};
\end{tikzpicture}
\caption{In the laboratory the \(z\)-directed force is magnetic; in the
instantaneous co-moving frame the same force must be electric.}
\end{figure}

Now introduce an observer moving with the particle at that instant. The
particle's velocity in the new frame is zero, so the magnetic term
\(q\boldsymbol v'\times\boldsymbol B'\) vanishes. Yet in the slow-motion
limit both observers must agree that the particle has a \(z\)-acceleration.
If the Lorentz force law has the same form in both frames, the co-moving
observer must attribute the force to an electric field \(\boldsymbol E'\)
along \(z\). A field that was purely magnetic in the laboratory is therefore
not purely magnetic in the moving frame.

This is the physical reason to seek one object whose components include both
fields. A change of frame does not merely rotate \(\boldsymbol E\) among its
own three components and \(\boldsymbol B\) among its own three components; it
mixes the two sets.

\section{The electromagnetic field tensor}

Together the electric and magnetic fields have six components. A
four-dimensional antisymmetric tensor also has six independent components,
since
\begin{equation}
  \frac{4(4-1)}{2}=6.
\end{equation}
This does not mean that an antisymmetric tensor is literally the only
collection of six quantities one could invent: six scalars, or a vector and
two scalars, would also contain six numbers. It is the natural geometrical
object whose components can mix in the required way.

To keep the board convention explicit, we display the contravariant tensor
in the spatial-first order \((1,2,3,0)\):
\begin{equation}
  F^{\mu\nu}=
  \begin{pmatrix}
    0 & -B_3 & B_2 & E_1 \\
    B_3 & 0 & -B_1 & E_2 \\
    -B_2 & B_1 & 0 & E_3 \\
    -E_1 & -E_2 & -E_3 & 0
  \end{pmatrix}.
  \label{eq:l07-field-tensor}
\end{equation}
The upper-left \(3\times3\) block is precisely the magnetic tensor
\eqref{eq:magnetic-tensor}; components with one spatial and one time index
contain the electric field. Antisymmetry supplies the remaining entries.

\begin{figure}[H]
\centering
\includegraphics[width=0.72\textwidth]{lecture_07_figure_08.png}
\caption{The magnetic spatial block and electric time-space column being
assembled into the four-dimensional field tensor.}
\lectureframe{00:46:35}
\end{figure}

\begin{classroomqa}
\audiencequestion{When should these field components carry upper rather
than lower indices?}
\lecturerresponse{For the three spatial labels alone, the board sometimes
suppresses the distinction while identifying the familiar Maxwell--Faraday
fields. In a covariant calculation the distinction must be restored: raising
or lowering a spatial index changes a sign in the chosen metric convention.}
\end{classroomqa}

What is invariant is not a particular arrangement of signs on the board,
but the antisymmetric spacetime tensor together with a consistent metric and
index convention. The notation \(F\) may suggest Faraday, although the
lecture does not claim that historical attribution with confidence. Einstein
worked out the transformation of electric and magnetic components in 1905;
the compact four-dimensional tensor language came with the later spacetime
formulation associated with Minkowski.

\section{Four-velocity and four-force}

The particle side of the equation must also be made covariant. Proper time
\(\tau\) is invariant, so differentiating the position four-vector with
respect to it gives another four-vector:
\begin{equation}
  u^\mu=\frac{dx^\mu}{d\tau}
  =\gamma(v_x,v_y,v_z,1),
  \qquad
  \gamma=\frac{1}{\sqrt{1-v^2}}.
  \label{eq:four-velocity}
\end{equation}
The last component follows from \(dt/d\tau=\gamma\). The four-momentum is
\begin{equation}
  p^\mu=mu^\mu,
\end{equation}
whose spatial part is the relativistic momentum
\(\boldsymbol p=m\gamma\boldsymbol v\).

\begin{figure}[H]
\centering
\includegraphics[width=0.75\textwidth]{lecture_07_figure_09.png}
\caption{Four-velocity, its components, the Lorentz factor, and
four-momentum.}
\lectureframe{00:50:30}
\end{figure}

Because \(p^\mu\) is a four-vector and \(\tau\) a scalar, the natural
relativistic left-hand side is
\begin{equation}
  f^\mu\equiv\frac{dp^\mu}{d\tau}.
  \label{eq:four-force-def}
\end{equation}
The four-force cannot be chosen arbitrarily. Four-velocity has fixed norm,
\begin{equation}
  u_\mu u^\mu=1.
\end{equation}
Differentiating along the world line gives
\begin{align}
  0
  &=\frac{d}{d\tau}(u_\mu u^\mu) \\
  &=\frac{du_\mu}{d\tau}u^\mu
    +u_\mu\frac{du^\mu}{d\tau} \\
  &=2u_\mu\frac{du^\mu}{d\tau}.
\end{align}
For constant mass this implies
\begin{equation}
  f^\mu u_\mu=0.
  \label{eq:force-orthogonal}
\end{equation}
Thus four-force is Minkowski-orthogonal to four-velocity. This is a purely
kinematic constraint following from the normalization of \(u^\mu\), not a
special property of electromagnetism.

\begin{figure}[H]
\centering
\includegraphics[width=0.64\textwidth]{lecture_07_figure_10.png}
\caption{The unit norm of four-velocity and the resulting orthogonality of
four-force and four-velocity.}
\lectureframe{00:56:50}
\end{figure}

The three-dimensional magnetic calculation has already supplied the trick
for satisfying \eqref{eq:force-orthogonal}. Contract an antisymmetric tensor
with one copy of the velocity:
\begin{equation}
  \boxed{
  \frac{dp^\mu}{d\tau}=qF^{\mu\nu}u_\nu.}
  \label{eq:covariant-lorentz}
\end{equation}
Indeed,
\begin{equation}
  f^\mu u_\mu
  =qF^{\mu\nu}u_\nu u_\mu=0,
\end{equation}
because \(F^{\mu\nu}\) is antisymmetric and \(u_\nu u_\mu\) is symmetric.
This is not the only conceivable relativistic force law, but it is the
simplest one linear in both the electromagnetic field and the velocity.

\begin{figure}[H]
\centering
\includegraphics[width=0.72\textwidth]{lecture_07_figure_11.png}
\caption{The covariant Lorentz-force equation on the left and its vanishing
contraction with a second four-velocity on the right.}
\lectureframe{01:00:00}
\end{figure}

\section{Recovering the ordinary Lorentz force}

A covariant-looking equation must still reproduce the known force law. With
the metric convention stated at the beginning,
\begin{equation}
  u_\mu=\gamma(-v_x,-v_y,-v_z,1).
\end{equation}
For the \(x\)-component, \eqref{eq:l07-field-tensor} and
\eqref{eq:covariant-lorentz} give
\begin{align}
  \frac{dp^1}{d\tau}
  &=q\bigl(F^{12}u_2+F^{13}u_3+F^{10}u_0\bigr) \\
  &=q\gamma\bigl(B_3v_y-B_2v_z+E_1\bigr) \\
  &=q\gamma\bigl[\boldsymbol E+\boldsymbol v\times
    \boldsymbol B\bigr]_x.
  \label{eq:x-force}
\end{align}
The other spatial components follow by cycling the indices.

\begin{classroomqa}
\audiencequestion{Where is the \(qE_x\) term in the component expansion?}
\lecturerresponse{It is the term \(qF^{10}u_0\). A field-tensor component
with one spatial and one time index is an electric field component, and
\(u_0=\gamma\).}
\end{classroomqa}

The factor \(\gamma\) on the right is not an extra force. Since
\begin{equation}
  \frac{dp^i}{d\tau}=\frac{dt}{d\tau}\frac{dp^i}{dt}
  =\gamma\frac{dp^i}{dt},
\end{equation}
it cancels from both sides of \eqref{eq:x-force}. The exact spatial equation
is therefore again
\begin{equation}
  \frac{d\boldsymbol p}{dt}
  =q\bigl(\boldsymbol E+\boldsymbol v\times\boldsymbol B\bigr),
  \qquad
  \boldsymbol p=m\gamma\boldsymbol v.
  \label{eq:exact-three-force}
\end{equation}
The familiar form survives, but the relativity is hidden in the momentum.

\begin{figure}[H]
\centering
\includegraphics[width=0.75\textwidth]{lecture_07_figure_12.png}
\caption{The three-dimensional Lorentz law recovered after the common
proper-time factor is removed; the momentum retains its factor of
\(\gamma\).}
\lectureframe{01:06:30}
\end{figure}

The time component carries the corresponding energy statement,
\begin{equation}
  \frac{dp^0}{dt}=q\boldsymbol E\cdot\boldsymbol v,
\end{equation}
so a magnetic field performs no work while an electric field may change the
particle energy.
\editorialnote{The time-component check is the direct companion of the
lecture's spatial calculation; it is included to expose the same no-work
property in four-dimensional notation.}

\begin{classroomqa}
\audiencequestion{Is this derivation mathematics or physics?}
\lecturerresponse{It is both. Tensor algebra supplies the compact form, but
substantial physics entered through the principle that all inertial frames
obey the same laws and through the assumption that the equations of motion
are second order. A second-order law predicts the future from position and
velocity; a third-order law would require acceleration as additional initial
data. Experiment, not notation, tells us which kind of law nature uses.}
\end{classroomqa}

\section{How the fields transform}

Equation \eqref{eq:covariant-lorentz} has the same form in every frame because
both sides are four-vectors. The field tensor itself transforms with one
Lorentz matrix for each index:
\begin{equation}
  F'^{\mu\nu}
  =\Lambda^\mu{}_{\rho}\Lambda^\nu{}_{\sigma}F^{\rho\sigma}.
  \label{eq:tensor-transform}
\end{equation}
The board derives this rule by first transforming a simple product
\(A^\mu B^\nu\). Once the rule is known for a product of vectors, linearity
gives the rule for any rank-two tensor.

For a boost along the \(1\)-axis, the coordinates transform as
\begin{align}
  x'^1&=\gamma(x^1-vx^0), &
  x'^0&=\gamma(x^0-vx^1), \\
  x'^2&=x^2, & x'^3&=x^3.
\end{align}
Apply the same transformation separately to both tensor indices. For the
component parallel to the boost, the board begins with a product:
\begin{align}
  (A'^0B'^1)
  =\gamma^2\bigl(
     A^0B^1-vA^0B^0-vA^1B^1+v^2A^1B^0
   \bigr).
\end{align}
Replacing each product by the corresponding tensor component and then using
antisymmetry, \(F^{00}=F^{11}=0\) and \(F^{10}=-F^{01}\), gives
\begin{equation}
  F'^{01}
  =\gamma^2(1-v^2)F^{01}
  =F^{01}.
\end{equation}
Thus the electric field component parallel to the motion is unchanged. The
parallel magnetic component behaves similarly.

The transverse components are different. Following the primed-frame and
sign convention used on the board,
\begin{align}
  F'^{02}
  &=\gamma\bigl(F^{02}-vF^{12}\bigr), \\
  E'_2
  &=\gamma(E_2+vB_3).
  \label{eq:e-transform}
\end{align}
A transverse electric component acquires part of a magnetic component. The
remaining transformation laws follow from the same two-index rule and are
left as the lecture's exercise.

\begin{figure}[H]
\centering
\includegraphics[width=0.76\textwidth]{lecture_07_figure_13.png}
\caption{The transverse tensor transformation at the point where the
\(E'_2\) component is seen to contain a term proportional to \(vB_3\).}
\lectureframe{01:23:40}
\end{figure}

The sign in \eqref{eq:e-transform} depends on which frame is declared to move
with \(+v\), as well as on index and metric conventions. Its invariant
content does not: transverse electric and magnetic components mix, and the
covariant force equation remains valid after the transformation.

\begin{classroomqa}
\audiencequestion{For a current in a wire, can one move at some velocity
that makes the magnetic field vanish?}
\lecturerresponse{Moving with the conduction electrons is not enough because
the positive lattice charges then move in the new frame. A neutral
current-carrying wire is therefore subtler than a single moving lump of
charge. A purely magnetic field cannot be transformed completely away. If
electric and magnetic fields are both present, their transformed pieces can
cancel in special circumstances; for one uniformly moving charge, the
charge's rest frame is the simple frame in which its magnetic field
vanishes.}
\end{classroomqa}

The precise invariant test uses
\begin{equation}
  \boldsymbol B^2-\boldsymbol E^2,
  \qquad
  \boldsymbol E\cdot\boldsymbol B.
\end{equation}
A frame with \(\boldsymbol B'=0\) can exist only for an electric-dominated
field with \(\boldsymbol E\cdot\boldsymbol B=0\); a field that is purely
magnetic in one frame is magnetic-dominated and cannot meet that condition.
\editorialnote{These invariant criteria state precisely the qualification
reached in the classroom discussion; their derivation is deferred.}

\begin{classroomqa}
\audiencequestion{Do the electric scalar potential and magnetic vector
potential form a four-vector?}
\lecturerresponse{Yes. With convention-dependent signs they combine into the
four-potential \(A^\mu\). Maxwell's equations will make systematic use of
that object.}
\end{classroomqa}

\section{Derivatives of fields}

We now leave the trajectory of one particle and consider a field throughout
spacetime. Let
\begin{equation}
  \phi=\phi(x^1,x^2,x^3,x^0)
\end{equation}
be a scalar field. Its four derivatives form the covariant components of a
four-vector:
\begin{equation}
  \partial_\mu\phi
  \equiv\frac{\partial\phi}{\partial x^\mu}
  =\left(
    \frac{\partial\phi}{\partial x},
    \frac{\partial\phi}{\partial y},
    \frac{\partial\phi}{\partial z},
    \frac{\partial\phi}{\partial t}
  \right).
\end{equation}
The chain rule applied to the Lorentz coordinate transformation proves this
covariant transformation law. Raising the index reverses the spatial signs
but leaves the time sign unchanged:
\begin{equation}
  \partial^\mu\phi
  =\left(
    -\frac{\partial\phi}{\partial x},
    -\frac{\partial\phi}{\partial y},
    -\frac{\partial\phi}{\partial z},
    \frac{\partial\phi}{\partial t}
  \right).
\end{equation}

Every differentiation contributes a lower index. Thus
\(\partial_\nu V^\mu\) is a rank-two tensor with one upper and one lower
index, and differentiating a tensor produces an object with one additional
lower index. This is the basic rule needed to write field equations without
choosing a preferred inertial frame.

\section{The invariant scalar wave equation}

The simplest wave equations contain two derivatives. If the equation is to
have the same form in every inertial frame, its left-hand side should be a
scalar. Contracting the derivative indices gives the unique simplest choice:
\begin{equation}
  \boxed{\partial^\mu\partial_\mu\phi=0.}
  \label{eq:wave-covariant}
\end{equation}
Writing out the components in the present signature,
\begin{equation}
  \frac{\partial^2\phi}{\partial t^2}
  -\frac{\partial^2\phi}{\partial x^2}
  -\frac{\partial^2\phi}{\partial y^2}
  -\frac{\partial^2\phi}{\partial z^2}=0.
  \label{eq:wave-components}
\end{equation}
The contraction in \eqref{eq:wave-covariant} makes Lorentz invariance
manifest; equation \eqref{eq:wave-components} is the same statement written
component by component.

\begin{figure}[H]
\centering
\includegraphics[width=0.75\textwidth]{lecture_07_figure_14.png}
\caption{The contracted derivative form of the scalar wave equation and its
time-minus-space component expansion.}
\lectureframe{01:39:40}
\end{figure}

\begin{classroomqa}
\audiencequestion{Must the right-hand side of the wave equation be zero?}
\lecturerresponse{No. It may be any scalar compatible with the problem. A
constant times \(\phi\), or a nonlinear scalar function such as one
proportional to \(\phi^2\), still preserves covariance. Zero gives the
simplest free wave equation.}
\end{classroomqa}

The point of this scalar example is preparation. Maxwell's equations mix
electric and magnetic fields more intricately, but they too can be written so
that Lorentz invariance is visible from their index structure. That is the
problem to which we return next time.
